\scanpage{18}
On désigne par
\[
u_{pq}:\mathfrak S_p\times\mathfrak S_q\longrightarrow\mathfrak S_{p+q}
\]
l'injection canonique, et on pose, pour
\[
F\in\operatorname{Ob}\mathcal A(\mathfrak X,p),\quad G\in\operatorname{Ob}\mathcal A(\mathfrak X,q)
\]
\[
F*G=u_{pq*}F\boxtimes G\qquad\in\operatorname{Ob}\mathcal A(\mathfrak X,p+q)
\]
L'opération $F*G$ est ``comm'' et ``ass.'' dans un sens évident. \deleted{\unsure{On peut donc}}\inserted{\unsure{Ainsi la}} \deleted{définir} ``catégorie graduée'' $(\mathcal A(\mathfrak X,p))_{p\in\mathbb N}$ devient un ``Anneau gradué'' ; en plus les ``Structures d'anneaux'' sur les morceaux homogènes $\mathcal A(\mathfrak X,p)$ séparément.

\deleted{\unsure{Fonctorialité|Foncteurs}}
On définit d'autre part les
\[
T^n(F)=F^{\otimes n}\quad\inserted{(F\in\operatorname{Ob}\mathcal A\mathfrak X)}
\]
comme objets de $\mathcal A(\mathfrak X,n)$ ;
\deleted{plus généralement $F\in\operatorname{Ob}\mathcal A(\mathfrak X,m)$, alors $T^n(F)=F^{\otimes n}$ est un $\unsure{O}_{\mathfrak X}$-Module à groupe d'opérateurs $\mathfrak S_m\times\mathfrak S_n$, \inserted{\unsure{opère diagonalement}} et utilisant $\unsure{u_{m,n}}$ on obtient}
\margin{17}
